\chapter{Model Training}

All models were trained on the 408 training contracts (80\%) only. Training ran on a single Apple-Silicon laptop using the MPS (Metal) backend with \texttt{torch} 2.8 and \texttt{transformers} 5.x; the summarizer ran on CPU for the memory reasons explained in Section~\ref{sec:engineering}. Table~\ref{tab:hyper} collects the configurations.

\begin{table}[ht]
\centering
\caption{Training configurations for the three models.}
\label{tab:hyper}
\small
\begin{tabular}{lccc}
\toprule
 & \textbf{Presence (MIL)} & \textbf{Span (QA)} & \textbf{Summarizer}\\
\midrule
Backbone            & DistilBERT-base & DistilBERT-base & FLAN-T5-small\\
Task head           & Sequence classification & QA (start/end) & Seq2seq LM\\
Training examples   & 24{,}000 (of 53{,}662) & 8{,}000 (of 11{,}156) & 4{,}000 (of 11{,}156)\\
Max input length    & 256 tokens & 320 tokens & 256 in / 48 out\\
Epochs              & 2 & 2 & 2\\
Batch size          & 16 & 16 & 8\\
Learning rate       & $2\times10^{-5}$ & $3\times10^{-5}$ & $3\times10^{-4}$\\
Device              & MPS & MPS & CPU\\
Wall-clock          & $\approx$49--52 min & $\approx$23 min & $\approx$8 min\\
\bottomrule
\end{tabular}
\end{table}

\section{Optimization Schedule}
Table~\ref{tab:hyper} gives the peak learning rates, but a peak rate on its own does not describe how a run behaved --- the schedule that decays it does. All three models use the same optimizer configuration: AdamW with a \textbf{linear decay} schedule, \textbf{no warmup}, and gradient-norm clipping at 1.0. Figure~\ref{fig:lrsched} plots the three resulting schedules. The presence and span settings were read back from the \texttt{training\_args.bin} saved beside each checkpoint rather than from our notes, so what is plotted is what actually ran.

Two choices are worth defending. Skipping warmup is normally risky, but these are short fine-tuning runs (1{,}000--3{,}000 steps) starting from pre-trained weights at small learning rates, where the usual justification for warmup --- stabilizing early updates in a randomly initialized model --- barely applies; the smooth loss curves in Figures~\ref{fig:presloss} and~\ref{fig:sumloss} show no early instability. Clipping at a gradient norm of 1.0 was left on throughout, which matters more than it might seem on the MPS backend: it is the guard that keeps a single bad batch from producing the kind of divergence we hit with DeBERTa-v3 (Section~\ref{sec:engineering}).

\begin{figure}[ht]
\centering
\includegraphics[width=\textwidth]{lr_schedule.png}
\caption{Learning-rate schedule for each of the three training runs: linear decay to zero with no warmup phase, over two epochs. The dashed line marks the start of the second epoch. Note the different vertical scales --- the summarizer runs at $3\times10^{-4}$, an order of magnitude above the two DistilBERT runs, which is standard for T5-family models.}
\label{fig:lrsched}
\end{figure}

We should be straightforward about one limitation here. A per-step gradient-norm trace would be the natural companion to these curves, and it is the one diagnostic that would let a reader confirm the runs stayed numerically healthy rather than infer it from the loss. We did not retain it: the intermediate checkpoint directories were cleared after training, so only the final weights and the loss values quoted below survive. Re-running the three trainings with \texttt{logging\_steps} capturing \texttt{grad\_norm} would recover it in about ninety minutes of wall-clock time, and we note it as an omission rather than dress up what we have.

\section{Presence Classifier}
The TF--IDF + logistic-regression baseline trains in under a minute. On the 20\% test set it reaches micro-F1 0.775, weighted-F1 0.777, macro-F1 0.572 over all 41 categories, and 0.666 over the 34 categories that have at least ten test positives. The gap between micro and macro comes from the near-few-shot tail of the dataset, not from a modelling failure.

The window-level DistilBERT was trained on a stratified subsample of 24{,}000 window examples (keeping the positive rate at 39.2\%) so the run stays tractable on a laptop. Training and validation losses fell smoothly (0.320$\to$0.261 train, 0.278$\to$0.257 validation over two epochs; Figure~\ref{fig:presloss}), and the two curves stayed close together, so we saw no sign of overfitting at this scale.

\begin{figure}[ht]
\centering
\includegraphics[width=0.62\textwidth]{presence_loss.png}
\caption{Training and validation loss of the window-level presence classifier.}
\label{fig:presloss}
\end{figure}

\section{Span Extractor}
The QA model trained on 8{,}000 focused-window examples for two epochs. Before training we decode the token-labelled spans back into text for a few examples and check they match the gold answers --- SQuAD-style pipelines make it very easy to get the character-to-token offset conversion silently wrong, and this small sanity check catches that.

\section{Summarizer}
FLAN-T5-small fine-tuned on 4{,}000 clause--template pairs in roughly eight minutes on CPU. The loss collapsed quickly (0.141$\to$0.085 train, 0.096$\to$0.075 validation; Figure~\ref{fig:sumloss}). In hindsight this rapid collapse was the first hint of the template-memorization behaviour we analyse in Chapter 6 --- with only 41 unique targets, the task is simply easy to fit.

\begin{figure}[ht]
\centering
\includegraphics[width=0.62\textwidth]{summarizer_loss.png}
\caption{Fine-tuning loss of the FLAN-T5-small summarizer.}
\label{fig:sumloss}
\end{figure}

\section{Engineering Notes: Training on Consumer Hardware}
\label{sec:engineering}
A few practical findings may save time for anyone reproducing this work on Apple-Silicon machines:
\begin{itemize}[itemsep=2pt]
  \item \textbf{DeBERTa-v3 is unstable on MPS.} Our first-choice backbone produced NaN gradients on the MPS backend and was far too slow on CPU (about 89 seconds per step). We switched to DistilBERT; the method is unchanged, and on CUDA hardware the backbone string can simply be swapped back.
  \item \textbf{MPS memory does not release in-process.} After two transformer trainings in the same session, the summarizer hit the MPS memory ceiling even after we deleted objects and emptied the cache. The fixes that actually work are restarting the kernel or training the (small) summarizer on CPU, which is what we did.
  \item \textbf{API migration.} \texttt{transformers} 5.x renamed \texttt{Trainer(tokenizer=\dots)} to \texttt{processing\_class=\dots}. We also disabled mixed precision (\texttt{fp16=bf16=False}) and cast loss logits to fp32 for MPS dtype safety.
  \item \textbf{Reproducibility.} Every artifact --- the split (\texttt{artifacts/split.json}), the fitted baseline (\texttt{artifacts/baseline.pkl}), and the three model checkpoints --- is saved to disk, and a standalone evaluation notebook regenerates every number in Chapter 6 from those files alone.
\end{itemize}
